\chapter{Curvature}
\label{chap:dc-curvature}

\section{Curvature}

\begin{definition}[curvature]
  \label{def:dc-ch4-2-1}
  \lean{Riemannian.AffineConnection.curvature}
  \uses{def:dc-ch2-2-1}
  \leanok
  \dcref{ch4:2.1}
  The \emph{curvature} $R$ of a Riemannian manifold $M$ is a correspondence that
  associates to every pair $X,Y\in\mathcal{X}(M)$ a mapping
  $R(X,Y):\mathcal{X}(M)\to\mathcal{X}(M)$ given by
  $$
  R(X,Y)Z=\nabla_Y\nabla_X Z-\nabla_X\nabla_Y Z+\nabla_{[X,Y]}Z,\quad Z\in\mathcal{X}(M),
  $$
  where $\nabla$ is the Riemannian connection of $M$.
  For $M=\mathbb{R}^n$, the coordinate expression
  $\nabla_X Z=(Xz_1,\dots,Xz_n)$ gives $R(X,Y)Z=0$; thus $R$ measures the
  deviation from the Euclidean connection.
  In a coordinate system $\{x_i\}$ around $p\in M$, since
  $[\frac{\partial}{\partial x_i},\frac{\partial}{\partial x_j}]=0$,
  $$
  R\Big(\frac{\partial}{\partial x_i},\frac{\partial}{\partial x_j}\Big)\frac{\partial}{\partial x_k}=\Big(\nabla_{\partial/\partial x_j}\nabla_{\partial/\partial x_i}-\nabla_{\partial/\partial x_i}\nabla_{\partial/\partial x_j}\Big)\frac{\partial}{\partial x_k},
  $$
  Hence curvature measures the non-commutativity of covariant derivatives. We use
  this sign convention; the opposite sign also occurs in the literature.
\end{definition}

\begin{proposition}[linearity properties of $R$]
  \label{prop:dc-ch4-2-2}
  \lean{Riemannian.AffineConnection.curvature_add_left, Riemannian.AffineConnection.curvature_add_middle, Riemannian.AffineConnection.curvature_add_right, Riemannian.AffineConnection.curvature_smul_left, Riemannian.AffineConnection.curvature_smul_middle, Riemannian.AffineConnection.curvature_smul_right}
  \uses{def:dc-ch4-2-1}
  \leanok
  \dcref{ch4:2.2}
  The curvature $R$ has the following properties:
  \begin{enumerate}
  \item[(i)] $R$ is bilinear in $\mathcal{X}(M)\times\mathcal{X}(M)$:
  $$
  R(fX_1+gX_2,Y_1)=fR(X_1,Y_1)+gR(X_2,Y_1),
  $$
  $$
  R(X_1,fY_1+gY_2)=fR(X_1,Y_1)+gR(X_1,Y_2),\quad f,g\in\mathcal{D}(M),\ X_1,X_2,Y_1,Y_2\in\mathcal{X}(M).
  $$
  \item[(ii)] For $X,Y\in\mathcal{X}(M)$, the operator
  $R(X,Y):\mathcal{X}(M)\to\mathcal{X}(M)$ is linear, that is,
  $$
  R(X,Y)(Z+W)=R(X,Y)Z+R(X,Y)W,\qquad R(X,Y)fZ=fR(X,Y)Z,
  \quad f\in\mathcal{D}(M),\ Z,W\in\mathcal{X}(M).
  $$
  \end{enumerate}
\end{proposition}
\begin{proof}
  \leanok
  \sketch
  Bilinearity in (i) follows by expanding the definition and applying the
  connection axioms in the first two arguments. Additivity in (ii) is immediate.
  For homogeneity,
  
  $$
  \nabla_Y\nabla_X(fZ)=\nabla_Y(f\nabla_X Z+(Xf)Z)=f\nabla_Y\nabla_X Z+(Yf)(\nabla_X Z)+(Xf)(\nabla_Y Z)+(Y(Xf))Z.
  $$
  
  Therefore
  
  $$
  \nabla_Y\nabla_X(fZ)-\nabla_X\nabla_Y(fZ)=f(\nabla_Y\nabla_X-\nabla_X\nabla_Y)Z+((YX-XY)f)Z,
  $$
  
  hence
  
  $$
  R(X,Y)fZ=f\nabla_Y\nabla_X Z-f\nabla_X\nabla_Y Z+([Y,X]f)Z+f\nabla_{[X,Y]}Z+([X,Y]f)Z=fR(X,Y)Z.\ \blacksquare
  $$
\end{proof}

\begin{remark}[Remark 2.3]
  \label{rem:dc-ch4-2-3}
  \uses{rem:dc-ch4-2-6}
  \dcref{ch4:2.3}
  The term $\nabla_{[X,Y]}Z$ is exactly what makes $R(X,Y)$ linear in its
  third argument; the coordinate form is recorded in \cref{rem:dc-ch4-2-6}.
\end{remark}

\begin{proposition}[Bianchi Identity]
  \label{prop:dc-ch4-2-4}
  \lean{Riemannian.AffineConnection.curvature_bianchi}
  \uses{def:dc-ch4-2-1, def:dc-ch2-3-4, lem:dc-ch0-5-3-jacobi}
  \leanok
  \dcref{ch4:2.4}
  $$
  R(X,Y)Z+R(Y,Z)X+R(Z,X)Y=0.
  $$
\end{proposition}
\begin{proof}
  \leanok
  \sketch
  Expanding the definition and using torsion-freeness gives
  
  $$
  \begin{aligned}R(X,Y)Z+R(Y,Z)X+R(Z,X)Y&=\nabla_Y\nabla_X Z-\nabla_X\nabla_Y Z+\nabla_{[X,Y]}Z\\&\quad+\nabla_Z\nabla_Y X-\nabla_Y\nabla_Z X+\nabla_{[Y,Z]}X\\&\quad+\nabla_X\nabla_Z Y-\nabla_Z\nabla_X Y+\nabla_{[Z,X]}Y\\&=\nabla_Y[X,Z]+\nabla_Z[Y,X]+\nabla_X[Z,Y]-\nabla_{[X,Z]}Y-\nabla_{[Y,X]}Z-\nabla_{[Z,Y]}X\\&=[Y,[X,Z]]+[Z,[Y,X]]+[X,[Z,Y]]=0,\end{aligned}
  $$
  
  where the last equality follows from the Jacobi identity for vector fields.
  $\blacksquare$
\end{proof}

\begin{proposition}[symmetries of the curvature]
  \label{prop:dc-ch4-2-5}
  \lean{Riemannian.AffineConnection.curvatureForm_antisymm_left, Riemannian.AffineConnection.curvatureForm_antisymm_right, Riemannian.AffineConnection.curvatureForm_bianchi, Riemannian.AffineConnection.curvatureForm_pairSwap}
  \uses{def:dc-ch4-2-1, prop:dc-ch4-2-4, def:dc-ch2-3-4, cor:dc-ch2-3-3}
  \leanok
  \dcref{ch4:2.5}
  % (a) = curvatureForm_bianchi (= prop:dc-ch4-2-4); (b) = curvatureForm_antisymm_left;
  % (c) = curvatureForm_antisymm_right (metric compatibility); (d) = curvatureForm_pairSwap.
  Write $\langle R(X,Y)Z,T\rangle=(X,Y,Z,T)$. Then
  \begin{enumerate}
  \item[(a)] $(X,Y,Z,T)+(Y,Z,X,T)+(Z,X,Y,T)=0$;
  \item[(b)] $(X,Y,Z,T)=-(Y,X,Z,T)$;
  \item[(c)] $(X,Y,Z,T)=-(X,Y,T,Z)$;
  \item[(d)] $(X,Y,Z,T)=(Z,T,X,Y)$.
  \end{enumerate}
\end{proposition}
\begin{proof}
  \leanok
  \sketch
  (a) is just the Bianchi identity again; (b) follows directly from \cref{def:dc-ch4-2-1};
  (c) is equivalent to $(X,Y,Z,Z)=0$, whose proof follows:
  
  $$
  (X,Y,Z,Z)=\langle\nabla_Y\nabla_X Z-\nabla_X\nabla_Y Z+\nabla_{[X,Y]}Z,Z\rangle.
  $$
  
  But
  
  $$
  \langle\nabla_Y\nabla_X Z,Z\rangle=Y\langle\nabla_X Z,Z\rangle-\langle\nabla_X Z,\nabla_Y Z\rangle,\qquad\langle\nabla_{[X,Y]}Z,Z\rangle=\tfrac{1}{2}[X,Y]\langle Z,Z\rangle.
  $$
  
  Hence
  
  $$
  \begin{aligned}(X,Y,Z,Z)&=Y\langle\nabla_X Z,Z\rangle-X\langle\nabla_Y Z,Z\rangle+\tfrac{1}{2}[X,Y]\langle Z,Z\rangle\\&=\tfrac{1}{2}Y(X\langle Z,Z\rangle)-\tfrac{1}{2}X(Y\langle Z,Z\rangle)+\tfrac{1}{2}[X,Y]\langle Z,Z\rangle\\&=-\tfrac{1}{2}[X,Y]\langle Z,Z\rangle+\tfrac{1}{2}[X,Y]\langle Z,Z\rangle=0,\end{aligned}
  $$
  
  which proves (c). To prove (d), we use (a), and write:
  
  $$
  \begin{aligned}(X,Y,Z,T)+(Y,Z,X,T)+(Z,X,Y,T)&=0,\\(Y,Z,T,X)+(Z,T,Y,X)+(T,Y,Z,X)&=0,\\(Z,T,X,Y)+(T,X,Z,Y)+(X,Z,T,Y)&=0,\\(T,X,Y,Z)+(X,Y,T,Z)+(Y,T,X,Z)&=0.\end{aligned}
  $$
  
  Summing the equations above we obtain $2(Z,X,Y,T)+2(T,Y,Z,X)=0$ and therefore
  $(Z,X,Y,T)=(Y,T,Z,X)$. $\blacksquare$
\end{proof}

\begin{remark}[Remark 2.6]
  \label{rem:dc-ch4-2-6}
  \lean{Riemannian.leviCivita_curvature_chartFrame_expansion, Riemannian.AffineConnection.curvature_apply_congr, Riemannian.AffineConnection.curvatureFormAt, Riemannian.AffineConnection.isAlgCurvatureForm_curvatureFormAt, Riemannian.leviCivita_curvatureFormAt_chartFrame}
  \uses{rem:dc-ch2-2-3, def:dc-ch2-3-7-christoffel, prop:dc-ch4-2-2, prop:dc-ch4-2-5}
  \leanok
  \dcref{ch4:2.6}
  It is convenient to express what was seen above in a coordinate system $(U,\mathbf{x})$
  based at $p\in M$. Indicating $\frac{\partial}{\partial x_i}=X_i$, we put
  $$
  R(X_i,X_j)X_k=\sum_\ell R_{ijk}^\ell X_\ell,
  $$
  so the $R_{ijk}^\ell$ are the components of the curvature $R$ in $(U,\mathbf{x})$.
  If $X=\sum_i u^i X_i$, $Y=\sum_j v^j X_j$, $Z=\sum_k w^k X_k$, then by linearity of
  $R$,
  $$
  R(X,Y)Z=\sum_{i,j,k,\ell}R_{ijk}^\ell u^i v^j w^k X_\ell.\qquad(1)
  $$
  To express $R_{ijk}^\ell$ in terms of the coefficients $\Gamma_{ij}^k$ of the
  Riemannian connection, we write
  $$
  R(X_i,X_j)X_k=\nabla_{X_j}\nabla_{X_i}X_k-\nabla_{X_i}\nabla_{X_j}X_k=\nabla_{X_j}\Big(\sum_\ell\Gamma_{ik}^\ell X_\ell\Big)-\nabla_{X_i}\Big(\sum_\ell\Gamma_{jk}^\ell X_\ell\Big),
  $$
  which by a direct calculation yields
  $$
  R_{ijk}^s=\sum_\ell\Gamma_{ik}^\ell\Gamma_{j\ell}^s-\sum_\ell\Gamma_{jk}^\ell\Gamma_{i\ell}^s+\frac{\partial}{\partial x_j}\Gamma_{ik}^s-\frac{\partial}{\partial x_i}\Gamma_{jk}^s.\qquad(2)
  $$
  Putting $\langle R(X_i,X_j)X_k,X_s\rangle=\sum_\ell R_{ijk}^\ell g_{\ell s}=R_{ijks}$,
  the identities of Proposition 2.5 can be written as
  $$
  R_{ijks}+R_{jkis}+R_{kijs}=0,\qquad R_{ijks}=-R_{jiks},\qquad R_{ijks}=-R_{ijsk},\qquad R_{ijks}=R_{ksij}.
  $$
  Equation (1), which depends on the linearity of the operator $R$, shows that the
  value of $R(X,Y)Z$ at the point $p$ depends uniquely on the values of $X,Y,Z$ at
  $p$ and the values of the functions $R_{ijk}^\ell$ at $p$. This contrasts with the
  behavior of the covariant derivative (see \cref{rem:dc-ch2-2-3}), the reason
  being that the covariant derivative is not linear in all of its arguments. In
  general, entities such as the curvature that are linear are called \emph{tensors} on
  $M$.
\end{remark}

\section{Sectional curvature}

Given a vector space $V$, we denote by $|x\wedge y|$ the expression

$$
\sqrt{|x|^2|y|^2-\langle x,y\rangle^2},
$$

which represents the area of a two-dimensional parallelogram determined by the
pair of vectors $x,y\in V$.

\begin{proposition}[sectional curvature independent of basis]
  \label{prop:dc-ch4-3-1}
  \lean{Riemannian.IsAlgCurvatureForm.sectionalCurvature_changeBasis}
  \uses{prop:dc-ch4-2-5, def:dc-ch4-3-2}
  \leanok
  \dcref{ch4:3.1}
  Let $\sigma\subset T_pM$ be a two-dimensional subspace of the tangent space
  $T_pM$ and let $x,y\in\sigma$ be two linearly independent vectors. Then
  $$
  K(x,y)=\frac{(x,y,x,y)}{|x\wedge y|^2}
  $$
  does not depend on the choice of the vectors $x,y\in\sigma$.
\end{proposition}
\begin{proof}
  \leanok
  \sketch
  Any two bases of $\sigma$ are related by swaps, nonzero rescalings, and shears
  $\{x,y\}\mapsto\{x+\lambda y,y\}$. The curvature symmetries and the identity
  $|x\wedge y|^2=|x|^2|y|^2-\langle x,y\rangle^2$ show that the numerator and
  denominator transform by the same square factor, so $K(x,y)$ is invariant.
\end{proof}

\begin{definition}[sectional curvature]
  \label{def:dc-ch4-3-2}
  \lean{Riemannian.sectionalCurvature}
  \uses{def:dc-ch4-2-1}
  \leanok
  \dcref{ch4:3.2}
  Given a point $p\in M$ and a two-dimensional subspace $\sigma\subset T_pM$, the
  real number $K(x,y)=K(\sigma)$, where $\{x,y\}$ is any basis of $\sigma$, is
  called the \emph{sectional curvature} of $\sigma$ at $p$.
  The importance of the sectional curvature comes from the fact that knowledge of
  $K(\sigma)$ for all $\sigma$ determines the curvature $R$ completely. This is a
  purely algebraic fact:
\end{definition}

\begin{lemma}[sectional curvature determines $R$]
  \label{lem:dc-ch4-3-3}
  \lean{Riemannian.IsAlgCurvatureForm.ext}
  \uses{prop:dc-ch4-2-5}
  \leanok
  \dcref{ch4:3.3}
  Let $V$ be a vector space of dimension $\geq 2$, provided with an inner product
  $\langle\,,\,\rangle$. Let $R:V\times V\times V\to V$ and $R':V\times V\times V\to V$
  be tri-linear mappings such that conditions (a), (b), (c) and (d) of \cref{prop:dc-ch4-2-5} are satisfied by
  $$
  (x,y,z,t)=\langle R(x,y)z,t\rangle,\qquad (x,y,z,t)'=\langle R'(x,y)z,t\rangle.
  $$
  If $x,y$ are two linearly independent vectors, we may write
  $$
  K(\sigma)=\frac{(x,y,x,y)}{|x\wedge y|^2},\qquad K'(\sigma)=\frac{(x,y,x,y)'}{|x\wedge y|^2},
  $$
  where $\sigma$ is the bi-dimensional subspace generated by $x$ and $y$. If for all
  $\sigma\subset V$, $K(\sigma)=K'(\sigma)$, then $R=R'$.
\end{lemma}
\begin{proof}
  \leanok
  \sketch
  It suffices to prove that $(x,y,z,t)=(x,y,z,t)'$ for any $x,y,z,t\in V$.
  Observe first that, by hypothesis, $(x,y,x,y)=(x,y,x,y)'$ for all $x,y\in V$. Then
  
  $$
  (x+z,y,x+z,y)=(x+z,y,x+z,y)',
  $$
  
  hence
  
  $$
  (x,y,x,y)+2(x,y,z,y)+(z,y,z,y)=(x,y,x,y)'+2(x,y,z,y)'+(z,y,z,y)',
  $$
  
  and therefore $(x,y,z,y)=(x,y,z,y)'$ for all $x,y,z\in V$. Using what we have just
  proved,
  
  $$
  (x,y+t,z,y+t)=(x,y+t,z,y+t)',
  $$
  
  hence
  
  $$
  (x,y,z,t)+(x,t,z,y)=(x,y,z,t)'+(x,t,z,y)',
  $$
  
  which can be written as
  
  $$
  (x,y,z,t)-(x,y,z,t)'=(y,z,x,t)-(y,z,x,t)'.
  $$
  
  It follows that the expression $(x,y,z,t)-(x,y,z,t)'$ is invariant by cyclic
  permutations of the first three elements. Therefore, by (a) of \cref{prop:dc-ch4-2-5},
  
  $$
  3[(x,y,z,t)-(x,y,z,t)']=0,
  $$
  
  hence $(x,y,z,t)=(x,y,z,t)'$ for all $x,y,z,t\in V$. $\blacksquare$
\end{proof}

\begin{lemma}[constant curvature characterization]
  \label{lem:dc-ch4-3-4}
  \lean{Riemannian.IsAlgCurvatureForm.eq_smul_stdCurvForm_of_const, Riemannian.IsAlgCurvatureForm.const_of_eq_smul_stdCurvForm}
  \uses{lem:dc-ch4-3-3, prop:dc-ch4-2-5}
  \leanok
  \dcref{ch4:3.4}
  Let $M$ be a Riemannian manifold and $p$ a point of $M$. Define a tri-linear
  mapping $R':T_pM\times T_pM\times T_pM\to T_pM$ by
  $$
  \langle R'(X,Y,W),Z\rangle=\langle X,W\rangle\langle Y,Z\rangle-\langle Y,W\rangle\langle X,Z\rangle,
  $$
  for all $X,Y,W,Z\in T_pM$. Then $M$ has constant sectional curvature equal to
  $K_o$ if and only if $R=K_o R'$, where $R$ is the curvature of $M$.
\end{lemma}
\begin{proof}
  \leanok
  \sketch
  Assume that $K(p,\sigma)=K_o$ for all $\sigma\subset T_pM$, and set
  $\langle R'(X,Y,W),Z\rangle=(X,Y,W,Z)'$. Observe that $R'$ satisfies the properties
  (a), (b), (c) and (d) of \cref{prop:dc-ch4-2-5}. Since
  
  $$
  (X,Y,X,Y)'=\langle X,X\rangle\langle Y,Y\rangle-\langle X,Y\rangle^2,
  $$
  
  we have that, for all pairs of vectors $X,Y\in T_pM$,
  
  $$
  R(X,Y,X,Y)=K_o(|X|^2|Y|^2-\langle X,Y\rangle^2)=K_o R'(X,Y,X,Y).
  $$
  
  \cref{lem:dc-ch4-3-3} implies that, for all $X,Y,W,Z$,
  
  $$
  R(X,Y,W,Z)=K_o R'(X,Y,W,Z),
  $$
  
  hence $R=K_o R'$. The converse is immediate. $\blacksquare$
\end{proof}

\begin{corollary}[constant curvature in an orthonormal basis]
  \label{cor:dc-ch4-3-5}
  \lean{Riemannian.IsAlgCurvatureForm.eq_kronecker_iff_const}
  \uses{lem:dc-ch4-3-4, prop:dc-ch4-2-5}
  \leanok
  \dcref{ch4:3.5}
  Let $M$ be a Riemannian manifold, $p$ a point of $M$ and $\{e_1,\dots,e_n\}$,
  $n=\dim M$, an orthonormal basis of $T_pM$. Define
  $R_{ijk\ell}=\langle R(e_i,e_j)e_k,e_\ell\rangle$, $i,j,k,\ell=1,\dots,n$. Then
  $K(p,\sigma)=K_o$ for all $\sigma\subset T_pM$ if and only if
  $$
  R_{ijk\ell}=K_o(\delta_{ik}\delta_{j\ell}-\delta_{i\ell}\delta_{jk}),
  $$
  where $\delta_{ij}=1$ if $i=j$ and $\delta_{ij}=0$ if $i\neq j$; equivalently,
  $K(p,\sigma)=K_o$ for all $\sigma\subset T_pM$ if and only if
  $R_{ijij}=-R_{ijji}=K_o$ for all $i\neq j$, and $R_{ijk\ell}=0$ in the other
  cases.
\end{corollary}
\begin{proof}
  \leanok
  \uses{lem:dc-ch4-3-4, lem:dc-ch4-3-3}
  \sketch
  By \cref{lem:dc-ch4-3-4}, $K(p,\sigma)=K_o$ for all $\sigma$ is equivalent to
  $R=K_oR'$, where $R'(X,Y,Z,T)=\langle X,Z\rangle\langle Y,T\rangle-\langle
  Y,Z\rangle\langle X,T\rangle$. Evaluating $R'$ on the orthonormal basis gives
  $R'(e_i,e_j,e_k,e_\ell)=\delta_{ik}\delta_{j\ell}-\delta_{i\ell}\delta_{jk}$,
  so $R=K_oR'$ yields
  $R_{ijk\ell}=K_o(\delta_{ik}\delta_{j\ell}-\delta_{i\ell}\delta_{jk})$.
  Conversely, since an algebraic curvature form is determined by its components
  in a basis (multilinearity, \cref{lem:dc-ch4-3-3}), the Kronecker identity on
  all basis $4$-tuples forces $R=K_oR'$, hence constant sectional curvature
  $K_o$. $\blacksquare$
\end{proof}

\section{Ricci curvature and scalar curvature}

For a unit vector $x=z_n\in T_pM$ and an orthonormal basis
$\{z_1,\dots,z_{n-1}\}$ of $x^\perp$, set

$$
\mathrm{Ric}_p(x)=\frac{1}{n-1}\sum_i\langle R(x,z_i)x,z_i\rangle,\quad i=1,2,\dots,n-1,
$$

$$
K(p)=\frac{1}{n}\sum_j\mathrm{Ric}_p(z_j)=\frac{1}{n(n-1)}\sum_{ij}\langle R(z_i,z_j)z_i,z_j\rangle,\quad j=1,\dots,n.
$$

These are the Ricci curvature in direction $x$ and the normalized scalar
curvature. Define the Ricci form by

$$
Q(x,y)=\text{trace of the mapping }z\mapsto R(x,z)y.
$$

For $x$ unit and $\{z_1,\dots,z_{n-1},z_n=x\}$ orthonormal,

$$
Q(x,y)=\sum_i\langle R(x,z_i)y,z_i\rangle=\sum_i\langle R(y,z_i)x,z_i\rangle=Q(y,x),
$$

Thus $Q$ is symmetric and $Q(x,x)=(n-1)\mathrm{Ric}_p(x)$, so Ricci curvature is
independent of the chosen orthonormal basis. If $K$ is defined by
$\langle K(x),y\rangle=Q(x,y)$, then for any orthonormal basis
$\{z_1,\dots,z_n\}$,

$$
\text{Trace of }K=\sum_j\langle K(z_j),z_j\rangle=\sum_j Q(z_j,z_j)=(n-1)\sum_j\mathrm{Ric}_p(z_j)=n(n-1)K(p),
$$

Consequently the normalized scalar curvature is basis-independent. The bilinear
form $\frac{1}{n-1}Q$ is the Ricci tensor. In coordinates $(x_i)$ with
$X_i=\frac{\partial}{\partial x_i}$, $g_{ij}=\langle X_i,X_j\rangle$ and $g^{ij}$
the inverse matrix of $g_{ij}$, the coefficients of the bilinear form
$\frac{1}{n-1}Q$ in the basis $\{X_i\}$ are

$$
\frac{1}{n-1}R_{ik}=\frac{1}{n-1}\sum_j R_{ijk}^j=\frac{1}{n-1}\sum_{sj}R_{ijks}g^{sj},
$$

and

$$
K=\frac{1}{n(n-1)}\sum_{ik}R_{ik}g^{ik}.
$$

\begin{definition}[Ricci form]
  \label{def:dc-ch4-4-ricci-form}
  \lean{Riemannian.ricciForm, Riemannian.ricciForm_eq_sum}
  \uses{prop:dc-ch4-2-5}
  \leanok
  \dcref{ch4:4}
  Let $p\in M$ and let $B(x,y,z,t)=\langle R(x,y)z,t\rangle$ be the algebraic
  curvature form on $T_pM$. The \emph{Ricci form} is the bilinear form
  $$
  Q(x,y)=\text{trace of the mapping }z\mapsto R(x,z)y,
  $$
  the trace of the endomorphism of $T_pM$ whose associated bilinear form is
  $(z,w)\mapsto B(x,z,y,w)$. In any orthonormal basis $\{e_i\}$ of $T_pM$,
  $$
  Q(x,y)=\sum_i\langle R(x,e_i)y,e_i\rangle=\sum_i B(x,e_i,y,e_i),
  $$
  and this sum is independent of the orthonormal basis, since it is a trace. The
  \emph{Ricci curvature} in the direction of a unit vector $x$ is
  $\mathrm{Ric}_p(x)=\frac{1}{n-1}Q(x,x)$.
\end{definition}

\begin{proposition}[the Ricci form is symmetric]
  \label{prop:dc-ch4-4-ricci-symm}
  \lean{Riemannian.ricciForm_symm}
  \uses{def:dc-ch4-4-ricci-form, prop:dc-ch4-2-5}
  \leanok
  \dcref{ch4:4}
  The Ricci form is symmetric, $Q(x,y)=Q(y,x)$ for all $x,y\in T_pM$. In
  particular $\mathrm{Ric}_p(x)$ is intrinsically defined.
\end{proposition}
\begin{proof}
  \leanok
  \sketch
  In an orthonormal basis $\{e_i\}$ we have $Q(x,y)=\sum_i B(x,e_i,y,e_i)$. By the
  pair-swap symmetry (d) of \cref{prop:dc-ch4-2-5}, $B(x,e_i,y,e_i)=B(y,e_i,x,e_i)$,
  hence $Q(x,y)=\sum_i B(y,e_i,x,e_i)=Q(y,x)$. $\blacksquare$
\end{proof}

\begin{definition}[scalar curvature]
  \label{def:dc-ch4-4-scalar}
  \lean{Riemannian.scalarCurvature, Riemannian.scalarCurvature_eq_sum}
  \uses{def:dc-ch4-4-ricci-form}
  \leanok
  \dcref{ch4:4}
  The (unnormalized) \emph{scalar curvature} at $p$ is the trace of the Ricci
  form $Q$,
  $$
  \mathrm{trace}(Q)=\sum_j Q(e_j,e_j)=\sum_{ij}\langle R(e_i,e_j)e_i,e_j\rangle
    =\sum_{ij}B(e_i,e_j,e_i,e_j),
  $$
  again independent of the orthonormal basis $\{e_i\}$ (a trace of a trace). do
  Carmo's normalized scalar curvature is $K(p)=\frac{1}{n(n-1)}\mathrm{trace}(Q)$.
\end{definition}

\begin{lemma}[curvature on a parametrized surface]
  \label{lem:dc-ch4-4-1}
  \lean{Riemannian.Jacobi.surface_covariant_commutator, Riemannian.Jacobi.surface_covariant_commutator_of_eventually}
  \uses{def:dc-ch4-2-1, prop:dc-ch2-2-2}
  \leanok
  \dcref{ch4:4.1}
  Let $f:A\subset\mathbb{R}^2\to M$ be a parametrized surface with coordinates
  $(s,t)$, and let $V=V(s,t)$ be a vector field along $f$. Then
  $$
  \frac{D}{dt}\frac{D}{ds}V-\frac{D}{ds}\frac{D}{dt}V=R\Big(\frac{\partial f}{\partial s},\frac{\partial f}{\partial t}\Big)V.\qquad(3)
  $$
\end{lemma}
\begin{proof}
  \sketch
  In local coordinates
  $(U,\mathbf{x})$ based at $p\in M$. Let $V=\sum_i v^i X_i$, where $v^i=v^i(s,t)$ and
  $X_i=\frac{\partial}{\partial x_i}$. Then
  
  $$
  \frac{D}{ds}V=\sum_i v^i\frac{D}{ds}X_i+\sum_i\frac{\partial v^i}{\partial s}X_i,
  $$
  
  $$
  \frac{D}{dt}\Big(\frac{D}{ds}V\Big)=\sum_i v^i\frac{D}{dt}\frac{D}{ds}X_i+\sum_i\frac{\partial v^i}{\partial t}\frac{D}{ds}X_i+\sum_i\frac{\partial v^i}{\partial s}\frac{D}{dt}X_i+\sum_i\frac{\partial^2 v^i}{\partial t\partial s}X_i.
  $$
  
  Interchanging the roles of $s$ and $t$ and subtracting,
  
  $$
  \frac{D}{dt}\frac{D}{ds}V-\frac{D}{ds}\frac{D}{dt}V=\sum_i v^i\Big(\frac{D}{dt}\frac{D}{ds}X_i-\frac{D}{ds}\frac{D}{dt}X_i\Big).
  $$
  
  Putting $f(s,t)=(x_1(s,t),\dots,x_n(s,t))$, so that
  $\frac{\partial f}{\partial s}=\sum_j\frac{\partial x_j}{\partial s}X_j$ and
  $\frac{\partial f}{\partial t}=\sum_k\frac{\partial x_k}{\partial t}X_k$, we have
  $\frac{D}{ds}X_i=\sum_j\frac{\partial x_j}{\partial s}\nabla_{X_j}X_i$ and
  
  $$
  \frac{D}{dt}\frac{D}{ds}X_i=\sum_j\frac{\partial^2 x_j}{\partial t\partial s}\nabla_{X_j}X_i+\sum_{jk}\frac{\partial x_j}{\partial s}\frac{\partial x_k}{\partial t}\nabla_{X_k}\nabla_{X_j}X_i,
  $$
  
  or
  
  $$
  \Big(\frac{D}{dt}\frac{D}{ds}-\frac{D}{ds}\frac{D}{dt}\Big)X_i=\sum_{jk}\frac{\partial x_j}{\partial s}\frac{\partial x_k}{\partial t}(\nabla_{X_k}\nabla_{X_j}X_i-\nabla_{X_j}\nabla_{X_k}X_i).
  $$
  
  Joining everything together, we finally get
  
  $$
  \Big(\frac{D}{dt}\frac{D}{ds}-\frac{D}{ds}\frac{D}{dt}\Big)V=\sum_{ijk}v^i\frac{\partial x_j}{\partial s}\frac{\partial x_k}{\partial t}R(X_j,X_k)X_i=R\Big(\frac{\partial f}{\partial s},\frac{\partial f}{\partial t}\Big)V.\ \blacksquare
  $$
\end{proof}

\section{Tensors on Riemannian manifolds}

The space $\mathcal{X}(M)$ is a module over $\mathcal{D}(M)$.

\begin{definition}[tensor of order $r$]
  \label{def:dc-ch4-5-1}
  \lean{Riemannian.CovariantTensor2Predicate, Riemannian.CovariantTensor4Predicate}
  \leanok
  \dcref{ch4:5.1}
  A \emph{tensor} $T$ of order $r$ on a Riemannian manifold is a multilinear mapping
  $$
  T:\underbrace{\mathcal{X}(M)\times\cdots\times\mathcal{X}(M)}_{r\text{ factors}}\to\mathcal{D}(M).
  $$
  This means that given $Y_1,\dots,Y_r\in\mathcal{X}(M)$, $T(Y_1,\dots,Y_r)$ is a
  differentiable function on $M$, and that $T$ is linear in each argument, that is,
  $$
  T(Y_1,\dots,fX+gY,\dots,Y_r)=fT(Y_1,\dots,X,\dots,Y_r)+gT(Y_1,\dots,Y,\dots,Y_r),
  $$
  for all $X,Y\in\mathcal{X}(M)$, $f,g\in\mathcal{D}(M)$.
  A tensor is a pointwise object. Fix $p\in M$ and let $U$ be a neighborhood of $p$
  on which it is possible to define vector fields $E_1,\dots,E_n\in\mathcal{X}(U)$
  such that at each $q\in U$ the vectors $\{E_i(q)\}$, $i=1,\dots,n$, form a basis of
  $T_qM$; we say $\{E_i\}$ is a \emph{moving frame} on $U$. Writing
  $Y_1=\sum_{i_1}y_{i_1}E_{i_1},\dots,Y_r=\sum_{i_r}y_{i_r}E_{i_r}$, by linearity
  $$
  T(Y_1,\dots,Y_r)=\sum_{i_1,\dots,i_r}y_{i_1}\cdots y_{i_r}T(E_{i_1},\dots,E_{i_r}).
  $$
  The functions $T(E_{i_1},\dots,E_{i_r})=T_{i_1\dots i_r}$ on $U$ are called the
  \emph{components} of $T$ in the frame $\{E_i\}$. The value of $T(Y_1,\dots,Y_r)$ at $p$
  depends only on the values at $p$ of the components of $T$ and the values of
  $Y_1,\dots,Y_r$ at $p$.
\end{definition}

\begin{example}[curvature tensor]
  \label{ex:dc-ch4-5-2}
  \lean{Riemannian.AffineConnection.curvatureForm, Riemannian.AffineConnection.curvatureForm_isCovariantTensor4}
  \uses{def:dc-ch4-5-1, prop:dc-ch4-2-2, prop:dc-ch4-2-5}
  \leanok
  \dcref{ch4:5.2}
  The \emph{curvature tensor}

  $$
  R:\mathcal{X}(M)\times\mathcal{X}(M)\times\mathcal{X}(M)\times\mathcal{X}(M)\to\mathcal{D}(M)
  $$

  is defined by

  $$
  R(X,Y,Z,W)=\langle R(X,Y)Z,W\rangle,\quad X,Y,Z,W\in\mathcal{X}(M).
  $$

  The components in the coordinate frame $\{X_i=\frac{\partial}{\partial x_i}\}$
  are $R(X_i,X_j,X_k,X_\ell)=R_{ijk\ell}$.
\end{example}
\begin{proof}
  \leanok
  \sketch
  Multilinearity of $R(X,Y,Z,W)=\langle R(X,Y)Z,W\rangle$ over $\mathcal{D}(M)$
  in the first three arguments is the $\mathcal{D}(M)$-linearity of the curvature
  operator $R(X,Y)Z$ (\cref{prop:dc-ch4-2-2}: additivity and homogeneity
  $R(fX,Y)Z=R(X,fY)Z=R(X,Y)fZ=fR(X,Y)Z$); linearity in the fourth argument $W$
  is the bilinearity of the metric $\langle\cdot,\cdot\rangle$. $\blacksquare$
\end{proof}

\begin{example}[metric tensor]
  \label{ex:dc-ch4-5-3}
  \lean{Riemannian.metricTensor, Riemannian.metricTensor_isCovariantTensor2}
  \uses{def:dc-ch4-5-1}
  \leanok
  \dcref{ch4:5.3}
  The "\emph{metric tensor}" $G:\mathcal{X}(M)\times\mathcal{X}(M)\to\mathcal{D}(M)$ is
  defined by $G(X,Y)=\langle X,Y\rangle$, $X,Y\in\mathcal{X}(M)$. $G$ is a tensor of
  order 2 and its components in the frame $\{X_i\}$ are the coefficients $g_{ij}$ of
  the Riemannian metric in the given system of coordinates.
\end{example}
\begin{proof}
  \leanok
  \sketch
  $\mathcal{D}(M)$-bilinearity of $G(X,Y)=\langle X,Y\rangle$ is inherited slotwise
  from the bilinearity of the inner product $g_p$ on each tangent space:
  $G(fX+gY,Z)=fG(X,Z)+gG(Y,Z)$ and symmetrically in the second argument.
  $\blacksquare$
\end{proof}

\begin{example}[the connection is not a tensor]
  \label{ex:dc-ch4-5-4}
  \lean{Riemannian.AffineConnection.connectionForm, Riemannian.AffineConnection.connectionForm_not_isCovariantTensor}
  \uses{def:dc-ch4-5-1, def:dc-ch2-2-1}
  \leanok
  \dcref{ch4:5.4}
  The Riemannian connection $\nabla$ defined by

  $$
  \nabla:\mathcal{X}(M)\times\mathcal{X}(M)\times\mathcal{X}(M)\to\mathcal{D}(M),\qquad\nabla(X,Y,Z)=\langle\nabla_X Y,Z\rangle,\quad X,Y,Z\in\mathcal{X}(M),
  $$

  is \emph{not} a tensor, because $\nabla$ is not linear with respect to the argument
  $Y$.
\end{example}
\begin{proof}
  \leanok
  \sketch
  The form is $\mathcal{D}(M)$-linear in the direction slot $X$, by the
  $\mathcal{D}(M)$-linearity of $\nabla$ in its first argument, and in the metric
  slot $Z$, by bilinearity of $\langle\cdot,\cdot\rangle$. In the differentiated
  slot $Y$, however, the Leibniz rule of the connection (\cref{def:dc-ch2-2-1},
  property (iii)) gives
  $$
  \nabla(X, fY, Z)=f\,\nabla(X, Y, Z)+(Xf)\,\langle Y, Z\rangle .
  $$
  The correction term $(Xf)\langle Y,Z\rangle$ is not of the tensorial form
  $f\,\nabla(X,Y,Z)$: whenever it is nonzero at some point — which occurs on any
  manifold of positive dimension — the middle-slot homogeneity required of an
  order-$3$ tensor (\cref{def:dc-ch4-5-1}) fails, so $\nabla$ is not a tensor.
  $\blacksquare$
\end{proof}

\begin{remark}[Remark 5.5]
  \label{rem:dc-ch4-5-5}
  \lean{Riemannian.RiemannianMetric.metricToDualEquiv, Riemannian.RiemannianMetric.metricToDual_apply, Riemannian.RiemannianMetric.metricRiesz_inner}
  \uses{def:dc-ch1-2-1, def:dc-ch4-5-1}
  \leanok
  \dcref{ch4:5.5}
  On a differentiable manifold without a metric, covariant tensors and
  contravariant tensors, defined using the dual $\mathcal{X}^*(M)$, are distinct.
  A Riemannian metric associates to each $X\in\mathcal{X}(M)$ a unique
  $\omega\in\mathcal{X}^*(M)$ given by
  $$
  \omega(Y)=\langle X,Y\rangle,\quad\text{for all }Y\in\mathcal{X}(M).
  $$
  This correspondence identifies covariant and contravariant tensors.
\end{remark}

\begin{remark}[Remark 5.6]
  \label{rem:dc-ch4-5-6}
  \lean{Riemannian.AffineConnection.covectorTensor, Riemannian.AffineConnection.covectorTensor_apply}
  \uses{rem:dc-ch4-5-5, ex:dc-ch4-5-3}
  \leanok
  \dcref{ch4:5.6}
  Identify the field $X\in\mathcal{X}(M)$ with the tensor
  $X:\mathcal{X}(M)\to\mathcal{D}(M)$ given by $X(Y)=\langle X,Y\rangle$,
  for all $Y\in\mathcal{X}(M)$.
\end{remark}

\begin{definition}[covariant differential of a tensor]
  \label{def:dc-ch4-5-7}
  \lean{Riemannian.AffineConnection.covariantDifferential1, Riemannian.AffineConnection.covariantDifferential2}
  \uses{def:dc-ch4-5-1, def:dc-ch2-2-1}
  \leanok
  \dcref{ch4:5.7}
  Let $T$ be a tensor of order $r$. The \emph{covariant differential} $\nabla T$ of $T$
  is a tensor of order $(r+1)$ given by
  $$
  \nabla T(Y_1,\dots,Y_r,Z)=Z(T(Y_1,\dots,Y_r))-T(\nabla_Z Y_1,\dots,Y_r)-\cdots-T(Y_1,\dots,Y_{r-1},\nabla_Z Y_r).
  $$
  For each $Z\in\mathcal{X}(M)$, the \emph{covariant derivative} $\nabla_Z T$ of $T$
  relative to $Z$ is a tensor of order $r$ given by
  $$
  \nabla_Z T(Y_1,\dots,Y_r)=\nabla T(Y_1,\dots,Y_r,Z).
  $$
  In a convenient frame this becomes natural. Let $p\in M$ and let
  $\alpha:(-\varepsilon,\varepsilon)\to M$ be a differentiable curve with $\alpha(0)=p$,
  $\alpha'(t)=Z(\alpha(t))$. Let $\{e_1,\dots,e_n\}$ be a basis of $T_pM$ and let
  $e_i(t)$ be the parallel transport of $e_i$ along $\alpha$. Let $T_{i_1\dots i_r}(t)$
  be the components, in the basis $\{e_i(t)\}$, of the restriction $T(\alpha(t))$ of
  $T$ to the curve $\alpha$. Then, by the definition of $\nabla_Z T$,
  $$
  (\nabla_Z T)(e_{i_1}(t),\dots,e_{i_r}(t))=\frac{d}{dt}T_{i_1\dots i_r}(t)-T(\nabla_Z e_{i_1},\dots,e_{i_r}(t))-\cdots-T(e_{i_1}(t),\dots,\nabla_Z e_{i_r}(t)).
  $$
  Since $\nabla_Z e_i(t)=0$, we have by linearity
  $$
  (\nabla_Z T)_{i_1\dots i_r}=(\nabla_Z T)(e_{i_1}(t),\dots,e_{i_r}(t))=\frac{d}{dt}T_{i_1\dots i_r}.
  $$
  Thus, in this frame, the components of $\nabla_ZT$ are the ordinary derivatives
  of the components of $T$.
\end{definition}

\begin{example}[covariant differential of the metric tensor]
  \label{ex:dc-ch4-5-8}
  \lean{Riemannian.AffineConnection.covariantDifferential2_metricTensor_eq_zero}
  \uses{def:dc-ch4-5-7, ex:dc-ch4-5-3, def:dc-ch2-3-1}
  \leanok
  \dcref{ch4:5.8}
  The covariant differential of the metric tensor is the zero tensor. Indeed, for
  all $X,Y,Z\in\mathcal{X}(M)$,

  $$
  \nabla G(X,Y,Z)=Z\langle X,Y\rangle-\langle\nabla_Z X,Y\rangle-\langle X,\nabla_Z Y\rangle=0,
  $$

  because $\nabla$ is the Riemannian connection.
\end{example}
\begin{proof}
  \leanok
  \sketch
  By \cref{def:dc-ch4-5-7}, $\nabla G(X,Y,Z)=Z\langle X,Y\rangle-\langle\nabla_Z X,Y\rangle
  -\langle X,\nabla_Z Y\rangle$. Compatibility of the Levi-Civita connection with the
  metric (\cref{def:dc-ch2-3-1}) gives $Z\langle X,Y\rangle=\langle\nabla_Z X,Y\rangle
  +\langle X,\nabla_Z Y\rangle$, so the three terms cancel. $\blacksquare$
\end{proof}

\begin{example}[covariant derivative of a vector field as a tensor]
  \label{ex:dc-ch4-5-9}
  \lean{Riemannian.AffineConnection.covectorTensor, Riemannian.AffineConnection.covariantDifferential1_covectorTensor_eq}
  \uses{rem:dc-ch4-5-6, def:dc-ch4-5-7, def:dc-ch2-3-1}
  \leanok
  \dcref{ch4:5.9}
  Let $X\in\mathcal{X}(M)$. Identify $X$ with the tensor that associates to the
  vector field $Y\in\mathcal{X}(M)$ the function $\langle X,Y\rangle$ (see \cref{rem:dc-ch4-5-6}). The covariant derivative of the tensor $X$ relative to the vector field
  $Z\in\mathcal{X}(M)$ is such that, for all $Y\in\mathcal{X}(M)$,

  $$
  \nabla_Z X(Y)=\nabla X(Y,Z)=Z(X(Y))-X(\nabla_Z Y)=Z\langle X,Y\rangle-\langle X,\nabla_Z Y\rangle=\langle\nabla_Z X,Y\rangle.
  $$

  We conclude that the tensor $\nabla_Z X$ can be identified with the vector field
  $\nabla_Z X$. This justifies the notation adopted, and shows that the covariant
  derivative of tensors is a generalization of the covariant derivative of vector
  fields.
\end{example}
\begin{proof}
  \leanok
  \sketch
  By \cref{def:dc-ch4-5-7}, $\nabla_Z X(Y)=Z\langle X,Y\rangle-\langle X,\nabla_Z Y\rangle$.
  Metric compatibility (\cref{def:dc-ch2-3-1}) gives $Z\langle X,Y\rangle
  =\langle\nabla_Z X,Y\rangle+\langle X,\nabla_Z Y\rangle$, whence
  $\nabla_Z X(Y)=\langle\nabla_Z X,Y\rangle$; i.e. the tensor $\nabla_Z X$ is the
  covector tensor of the vector field $\nabla_Z X$. $\blacksquare$
\end{proof}
